\section{RECURSOS MATERIAIS E ORÇAMENTO PREVISTO}

Por se tratar de um projeto fundamentado nos princípios da Tecnologia Social e da sustentabilidade acadêmica, a arquitetura foi planejada para operar inteiramente sob camadas gratuitas (\textit{free tier}) de serviços de nuvem e ferramentas \textit{open source}, exigindo investimento financeiro direto de R\$ 0,00 para a instituição de ensino e para o pesquisador.

\begin{table}[!htb]
    \centering
    \caption{Recursos e Orçamento Previsto para o Projeto EloMaterno}
    \label{tab:orcamento}
    \small
    \begin{tabular}{p{4.2cm}p{6.5cm}p{3cm}r}
        \toprule
        \textbf{Recurso / Serviço} & \textbf{Finalidade Técnica} & \textbf{Provedor / Camada} & \textbf{Custo (R\$)} \\
        \midrule
        Estação de Desenvolvimento & Computador pessoal com Linux/Ubuntu, VS Code, Git e Docker & Infraestrutura própria & 0,00 \\
        \addlinespace
        Hospedagem \textit{Frontend} & Distribuição da SPA React em rede global de borda (CDN) com SSL & Vercel (Hobby Tier) & 0,00 \\
        \addlinespace
        Hospedagem \textit{Backend} API & Container Docker com JDK 17 slim executando o Spring Boot & Render (Free Web Service) & 0,00 \\
        \addlinespace
        Banco de Dados Gerenciado & Instância PostgreSQL 15+ gerenciada & Supabase / Neon (Free Tier) & 0,00 \\
        \addlinespace
        Armazenamento de Imagens & CDN de mídia com compressão dinâmica & Cloudinary (Free Tier) & 0,00 \\
        \addlinespace
        Repositório e CI/CD & Controle de versão e automação de testes com GitHub Actions & GitHub (Public Repo) & 0,00 \\
        \midrule
        \textbf{Total Geral} & \multicolumn{2}{l}{\textbf{Custo direto adicional para o projeto}} & \textbf{0,00} \\
        \bottomrule
    \end{tabular}
    \vspace{0.2cm}
    
    {\small Fonte: autor (2026).}
\end{table}